\subsubsection{Synchronisation des tuiles et choix du processeur}\label{sec:lancements-cooperatifs}
À 100\,000 bases, le calcul en bande demande encore plusieurs milliers de lancements. Le noyau coopératif conserve sur le GPU les deux lignes de scores et y lit la requête codée sur deux bits par base. Une barrière entre tous les blocs sépare deux tranches successives. Chaque lancement traite jusqu’à 256 tranches, soit 4\,096 lignes. La taille de la grille reste inférieure ou égale au nombre de blocs que le GPU peut exécuter simultanément.

La figure~\ref{fig:coop-focus} porte sur quinze paires de 100\,000 bases. Chaque méthode occupe une fois chacune des cinq positions dans l’ordre de passage. À 2\,\% d’éditions introduites, le temps moyen du score descend de 71,73 à 66,99\,ms avec les lancements coopératifs. Sur la première paire, Nsight Systems compte 25 lancements au lieu de 6\,249. La durée cumulée des noyaux passe de 65,24 à 66,20\,ms. Le regroupement raccourcit donc le temps écoulé sans réduire la durée cumulée observée sur le GPU.
\chart{comparison/nvpd_launch.png}{Comparaison des temps à 100\,000 bases sur GH200, rapportés à WFA2 sur un cœur Grace. À chaque taux, les cinq paires ont le même poids.}{fig:coop-focus}{bench/repeated/divergence/launch/analyze.py}

À 2\,\%, WFA2 prend 5,31\,ms. Un premier répartiteur choisit WFA2 dès qu’un alignement construit sur CPU coûte au plus \(\lfloor0{,}06\max(m,n)\rfloor\) éditions, pour des longueurs d’au plus 131\,072 bases. Décision comprise, les cinq paires prennent 5,51\,ms en moyenne. Après échauffement, le RSS hôte vaut 5,09\,Mio, contre 155,86\,Mio pour NvPD sur GPU. À 5\,\%, le chemin construit dépasse le seuil sur les cinq paires. Le répartiteur les envoie au GPU alors que WFA2 est plus rapide.

\synthese{À 100\,000 bases et 10\,\%, le noyau coopératif prend 68,90\,ms en moyenne, contre 100,18\,ms pour WFA2 sur un cœur. À 2\,\%, le rapport moyen entre les temps du noyau coopératif et de WFA2 vaut 12,63.}

\subsubsection{Choix du solveur à partir d’une borne d’alignement}
Un chemin glouton peut coûter plus de 6\,\% de la longueur alors que la distance optimale reste sous ce seuil. Le second répartiteur essaie alors de relier des mots communs de six bases par de courts alignements. Si le coût de l’un des deux chemins respecte le seuil, WFA2 calcule le score sur CPU. Sinon, les paires ne contenant que A, C, G et T, de 65\,536 à 131\,072 bases dont les longueurs diffèrent d’au plus un quart de la plus grande passent au noyau coopératif. Le seuil porte sur le coût d’un chemin construit. Le taux d’édition utilisé pour produire les séquences reste inconnu du répartiteur.

La figure~\ref{fig:dispatch-confirm} compare les deux politiques sur quinze nouvelles paires de 100\,000 bases. Dans chacun des six blocs, un processus lit les séquences, effectue un appel d’échauffement, puis chronomètre le score, choix du solveur compris. Les deux politiques utilisent WFA2 pour les cinq paires à 2\,\%. Le nouveau répartiteur prend 5,34\,ms en moyenne, soit 1,97\,\% de plus qu’un appel direct à WFA2.
\chart{comparison/nvpd_dispatch_confirm.png}{Temps par paire pour les deux politiques de choix du solveur et les appels directs, sur cinq paires distinctes à chaque taux. Les intervalles à 95\,\% portent sur les cinq moyennes par paire.}{fig:dispatch-confirm}{bench/repeated/divergence/launch/refine_figure.py}
À 5\,\%, la seconde borne oriente les cinq paires vers WFA2. Le temps moyen passe de 72,96 à 28,36\,ms. Comme le processus utilise uniquement le CPU, sa mémoire résidente proportionnelle après échauffement descend de 151,68 à 3,13\,Mio. À 10\,\%, les cinq paires passent au GPU. Le temps du répartiteur dépasse de 0,52\,\% celui du noyau coopératif appelé directement.

\synthese{À 5\,\%, la nouvelle politique réduit de 61,12\,\% le temps de l’ancienne (IC à 95\,\% de 60,86 à 61,39\,\%). À 10\,\%, elle réduit de 30,69\,\% le temps de WFA2 sur un cœur Grace (de 30,07 à 31,31\,\%).}
